\chapter{Background} \label{ch:background}

This chapter builds the conceptual foundation for covenant-based vault custody.
Section~\ref{se:script_limits} explains why Bitcoin Script cannot natively enforce
spending conditions on future transactions. Section~\ref{se:introspection} introduces
transaction introspection as the missing capability that covenants provide.
Section~\ref{se:covenant_taxonomy} distinguishes generic from purpose-built covenant
opcodes and discusses the recursive covenant question. Section~\ref{se:vault_lifecycle}
formalizes the vault lifecycle model and threat assumptions.
Sections~\ref{se:ctv}--\ref{se:catcsfs_bg} describe the four Bitcoin Script
extension proposals that are the primary subjects of this thesis.
Section~\ref{se:simplicity_bg} introduces Simplicity as an
\emph{alternative script primitive} on the Elements federated
sidechain, included as a cross-substrate reference point rather than a
peer of the four Bitcoin proposals.

% ======================================================================
\section{Bitcoin Script and Its Limitations} \label{se:script_limits}

Bitcoin~\cite{Nakamoto08} transactions transfer value between \emph{unspent transaction
outputs} (UTXOs)~\cite{EUTXO20, Antonopoulos17}. Each UTXO is encumbered by a
\emph{locking script} (scriptPubKey) that specifies the conditions under which the UTXO
may be spent. The spending transaction provides a \emph{witness} (scriptSig or Segwit
witness~\cite{BIP141}) that satisfies these conditions.

Bitcoin Script is deliberately limited. It is not Turing-complete: there are no loops,
no persistent state between transactions, and a hard 10{,}000-byte script size limit.
The language operates on a stack of byte strings and supports three categories of
operations relevant to custody:

\begin{itemize}
\item \textbf{Signature verification.} \texttt{OP\_CHECKSIG} and its variants verify
  that a provided digital signature matches a public key embedded in the script. This is
  the basis of all standard Bitcoin spending: proving possession of the private key
  corresponding to the output's public key.

\item \textbf{Hash preimage gates.} \texttt{OP\_HASH160}, \texttt{OP\_SHA256}, and
  related opcodes enable hash time-locked contracts (HTLCs), where spending requires
  revealing a preimage that hashes to a committed value. This is the foundation of the
  Lightning Network's payment channels.

\item \textbf{Timelocks.} \texttt{OP\_CHECKLOCKTIMEVERIFY} (CLTV) and
  \texttt{OP\_CHECKSEQUENCEVERIFY} (CSV)~\cite{BIP68, BIP112} enforce absolute and relative time constraints,
  respectively. A script can require that a UTXO remain unspent until a certain block
  height or for a certain number of blocks after its creation.
\end{itemize}

What Bitcoin Script \emph{cannot} do is constrain the \emph{structure} of the spending
transaction itself. A standard \texttt{OP\_CHECKSIG} script verifies that the spender
possesses a private key, but it places no restriction on where the funds go, how many
outputs the transaction creates, or what scripts encumber those outputs. Once the
signature check passes, the spender has full discretion over the transaction's outputs.

This limitation has a direct consequence for custody: a compromised private key grants
the attacker unconditional control over the funds. There is no mechanism within the
script to impose a time delay, require a second authorization, or force funds to a
predetermined recovery address. The custodian's only defense is to keep the key secure
--- there is no protocol-level fallback.

% ======================================================================
\section{Transaction Introspection} \label{se:introspection}

Transaction introspection is the ability of a script to examine fields of the
transaction that attempts to spend the UTXO it guards. Instead of merely verifying
a signature or hash preimage, an introspective script can inspect the spending
transaction's outputs, input count, sequence numbers, lock time, or other structural
fields and accept or reject the spend based on whether those fields match expected
values.

Introspection enables a qualitatively different class of spending conditions. Rather
than asking ``does the spender possess the right key?'' the script asks ``does the
spending transaction have the right structure?'' This distinction is the conceptual
foundation of covenants.

Consider a concrete example. Suppose Alice deposits 1~BTC into a UTXO whose script
enforces: ``this UTXO may only be spent by a transaction that creates exactly two
outputs---one sending 0.9~BTC to address~$X$ after a 144-block delay, and one
sending 0.1~BTC to address~$Y$ immediately.'' No standard Bitcoin Script can express
this condition, because the script has no access to the spending transaction's output
fields. An introspective script can: it reads the outputs at spend time and verifies
that they match the embedded constraints.

The degree of introspection varies across proposals:

\begin{itemize}
\item \textbf{Output-hash introspection.} The script computes a hash over all outputs
  of the spending transaction and compares it to an embedded commitment. CTV (BIP-119)
  and Simplicity's \texttt{jet::outputs\_hash()} use this approach. The script sees
  only the aggregate hash, not individual output fields.

\item \textbf{Per-output introspection.} The script examines individual outputs by
  index---verifying a specific output's scriptPubKey, amount, or both. CCV (BIP-443)
  operates at this granularity, checking that a particular output index matches an
  expected Taproot structure.

\item \textbf{Sighash preimage introspection.} The script reconstructs the transaction's
  signature hash from witness-provided fields and verifies it against both a stack-based
  signature check (\texttt{OP\_CHECKSIGFROMSTACK}) and the real transaction sighash
  (\texttt{OP\_CHECKSIG}). CAT+CSFS uses this technique to achieve introspection without
  a dedicated introspection opcode.
\end{itemize}

% ======================================================================
\section{Covenants: Generic and Purpose-Built} \label{se:covenant_taxonomy}

A \emph{covenant} is a spending condition that restricts not just who can spend a UTXO
but how the resulting transaction is structured~\cite{MES16}. The term originates from
property law, where deed covenants restrict how real estate may be used by future owners.
In Bitcoin, a covenant restricts how future transactions may spend a UTXO, propagating
constraints forward through the transaction graph.

\subsection{Generic Covenants} \label{ss:generic}

A \emph{generic covenant opcode} provides general-purpose introspection that can enforce
arbitrary constraints on transaction structure. OP\_CAT (BIP-347) is the minimal example:
by enabling stack concatenation, it allows scripts to assemble and verify sighash
preimages, achieving transaction introspection as a composition of
primitives~\cite{Poe21}. CTV (BIP-119) is generic in a different sense: it commits to a
complete transaction template, constraining the entire output structure in one hash
comparison.

Generic covenants are flexible---the same opcode can enforce vaults, payment channels,
congestion control, and other applications---but they impose the full complexity of
constraint specification on the script author. There are no guardrails: a subtle error
in the script (such as using an undefined mode value in CCV) can silently disable all
covenant enforcement.

\subsection{Purpose-Built Covenants} \label{ss:purpose_built}

A \emph{purpose-built covenant opcode} provides vault-specific semantics directly. OP\_VAULT
(BIP-345) is the primary example: the opcode pair \texttt{OP\_VAULT} and
\texttt{OP\_VAULT\_RECOVER} encodes the vault lifecycle (trigger withdrawal, authorize
recovery, enforce timelocks) at the opcode level. Amount preservation, partial
withdrawal, and recovery authorization are enforced at the opcode level rather than
composed from primitives.

Purpose-built covenants reduce the script author's burden and eliminate classes of
composition errors. The tradeoff is reduced generality: OP\_VAULT enforces only the
vault pattern, and extending it to other applications requires additional opcodes.

\subsection{The Recursive Covenant Question} \label{ss:recursive}

A covenant is \emph{recursive} if it can constrain its outputs to contain the same
covenant script, creating a chain of constrained transactions that extends indefinitely.
Recursive covenants enable patterns such as:

\begin{itemize}
\item \textbf{Perpetual vaults.} A vault that can be partially withdrawn from
  indefinitely, with the remainder re-locked under the same rules.
\item \textbf{State machines.} Multi-step protocols where each transaction advances the
  contract state and the next transaction is constrained by the current state.
\item \textbf{Congestion control trees.} A single UTXO that fans out into a tree of
  pre-committed transactions, allowing many recipients to claim funds at their own pace.
\end{itemize}

Whether a covenant opcode supports recursion depends on whether the script can reference
its own scriptPubKey in the output constraints. CTV does not support direct recursion:
the template hash commits to specific output scriptPubKeys, but the script cannot
compute its own hash at spend time. CCV supports recursion via its \emph{taptree
preservation} mode, which constrains an output to contain the same Taproot script
tree as the current input. OP\_VAULT supports a limited form through automatic revault
(excess value beyond the withdrawal amount is re-locked in a new vault UTXO).
CAT+CSFS can achieve one level of self-reference by embedding the script's own
\texttt{sha\_single\_output} as a constant, but cannot chain indefinitely: the
output script's hash would need to include itself as a constant, creating a
circular dependency that cannot be resolved at script construction
time. Simplicity
programs cannot reference their own commitment hash at runtime, placing them in the
same category as CTV for recursion purposes.

The recursive covenant question has practical consequences for vault design: it determines
whether partial withdrawals can be chained (relevant to the revault amplification and
per-UTXO recovery-cost scaling experiments in Chapters~\ref{ch:results} and~\ref{ch:security}).

% ======================================================================
\section{Vault Lifecycle and Threat Model} \label{se:vault_lifecycle}

The vault lifecycle follows the model formalized by Swambo et al.~\cite{SHMB20}. A
vault is a UTXO whose spending is constrained to a time-delayed withdrawal process
monitored by an automated watchtower. The lifecycle consists of four phases:

\begin{enumerate}
\item \textbf{Deposit.} The user sends funds to a vault address, creating a covenant-
  encumbered UTXO. The covenant enforces that this UTXO can only be spent via the
  trigger or recovery paths.

\item \textbf{Trigger (unvault).} The hot key initiates a withdrawal by broadcasting
  a trigger transaction. This moves funds from the vault UTXO to an unvaulting UTXO
  that is subject to a relative timelock (\texttt{OP\_CSV}).

\item \textbf{Withdrawal.} After the timelock expires (the ``watchtower window''),
  the hot key completes the withdrawal to the destination address. If the trigger was
  unauthorized, the watchtower intervenes during the delay.

\item \textbf{Recovery.} At any point before withdrawal completes, the cold key (or, in
  CCV's case, anyone) can invoke emergency recovery, sweeping funds to a pre-committed
  recovery address. Recovery bypasses the timelock---it is always immediately available.
\end{enumerate}

\begin{figure}[t]
\centering
\begin{tikzpicture}[
  node distance=3.5cm and 3.8cm,
  state/.style={rectangle, draw, rounded corners, minimum width=2.2cm,
                minimum height=0.9cm, font=\small},
  every edge/.style={draw, ->, >=stealth, thick},
]
  \node[state] (funded) {Funded};
  \node[state, right=of funded] (triggered) {Triggered};
  \node[state, right=of triggered] (withdrawn) {Withdrawn};
  \node[state, below=2cm of triggered] (recovered) {Recovered};

  \draw (funded) -- node[above, font=\footnotesize] {trigger (hot key)} (triggered);
  \draw (triggered) -- node[above, font=\footnotesize] {withdraw (hot + CSV)} (withdrawn);
  \draw (funded) -- node[left, font=\footnotesize, align=right, pos=0.3] {recover\\(cold key)} (recovered);
  \draw (triggered) -- node[right, font=\footnotesize] {recover} (recovered);
\end{tikzpicture}
\caption{Vault lifecycle state machine. All covenant designs implement this four-state
model. The trigger-to-withdrawal transition enforces a relative timelock (CSV) that
creates the watchtower monitoring window. Recovery is available from both the Funded
and Triggered states.}
\label{fig:vault_lifecycle}
\end{figure}

The watchtower is an automated process that monitors the blockchain for unauthorized
trigger transactions. When the watchtower detects an unexpected unvault, it broadcasts
a recovery transaction during the CSV delay window. The security of the vault depends
on two assumptions: (1)~the watchtower is online and connected to the network during
the delay period, and (2)~the recovery transaction can confirm before the attacker's
withdrawal.

\subsection{Threat Model Assumptions (Kerckhoffs)}
\label{ss:kerckhoffs}

We adopt the standard cryptographic threat model. The adversary is
assumed to know the vault's design, the covenant type, the script
structure, and all public parameters --- including the recovery
address where applicable. Security rests on keys and on the script's
semantics, not on the secrecy of construction. Every attack-feasibility
claim in this thesis is evaluated under this assumption. Concretely,
we do not rely on taproot leaf hiding, address non-disclosure, or
closed-source construction as security mechanisms.

This threat model choice is load-bearing for every attack-class
derivation in Chapter~\ref{ch:security}. A reviewer may object to
individual class-susceptibility claims with ``but taproot hides that
leaf,'' or ``but the attacker would not know the cold address''; those
objections target the on-chain information-reveal layer, where the
thesis makes no cryptographic claim. The Kerckhoffs assumption moves
the debate from ``does taproot leak?'' (which it does not, in the
absence of legitimate recovery use) to ``should a vault's griefing or
theft immunity depend on the secrecy of its construction?'' (which,
by Kerckhoffs, it should not). A vault that is griefing-immune only
because the attacker does not know its cold address is not
cryptographically secure; it is obscured. The moment the recovery path
is ever legitimately used on-chain, the cold address is revealed
permanently. Security-by-obscurity is therefore not a principled
defence, and the thesis assumes an attacker who has passed the
obscurity threshold.

\subsection{Threat Taxonomy} \label{ss:threat_taxonomy}

We extend the threat model of Swambo et al.~\cite{SHMB20} with attack classes
specific to the four covenant designs. Table~\ref{tab:threat_summary} summarizes the
eleven threat models used throughout this thesis. Each threat model specifies the
adversary's capabilities (which keys are compromised), the attack mechanism, and the
worst-case outcome. The full threat model matrix with per-covenant severity ratings
appears in Chapter~\ref{ch:security}.

\begin{table}[t]
\centering
\caption{Threat taxonomy. Each row defines an adversary profile used in the security
analysis (Chapter~\ref{ch:security}).}
\label{tab:threat_summary}
\small
\begin{tabular}{@{}clll@{}}
\toprule
ID & Attack class & Adversary capability & Worst-case outcome \\
\midrule
TM1  & Fee pinning            & Hot key + network access    & Fund theft (CTV) \\
TM2  & Recovery griefing      & Varies by covenant          & Denial of service \\
TM3  & Trigger key theft      & Hot/trigger key             & Fund theft or DoS \\
TM4  & Watchtower exhaustion  & Hot key + capital           & Fund loss (dust) \\
TM5  & Trigger xpub theft     & OP\_VAULT trigger xpub      & Amplified TM3 \\
TM6  & Dual-key compromise    & Hot + cold/auth key         & Fund theft or DoS \\
TM7  & Address reuse          & None (user error)           & Stuck funds (CTV) \\
TM8  & CCV mode bypass        & Script knowledge            & Total fund loss \\
TM9  & Hot key redirection    & Hot key (CAT+CSFS)          & Griefing only \\
TM10 & Witness tampering      & Hot key (CAT+CSFS)          & Impossible \\
TM11 & Cold key theft         & Cold key (CAT+CSFS)         & Total fund theft \\
\bottomrule
\end{tabular}
\end{table}

% ======================================================================
\section{OP\_CHECKTEMPLATEVERIFY (CTV)} \label{se:ctv}

CTV~\cite{BIP119} constrains the spending transaction to match a committed template
hash. The hash covers: \texttt{nVersion}, \texttt{nLockTime}, \texttt{scriptSigs} hash,
input count, \texttt{nSequence} hash, output count, outputs hash, and the current
input's index. A script of the form \texttt{<hash> OP\_CHECKTEMPLATEVERIFY} succeeds
only if the spending transaction's template hash matches exactly.

CTV commits to the transaction's output structure but \emph{not} to which UTXOs are
consumed as inputs (input prevouts and amounts are excluded from the hash). This
design choice means the same CTV script can be satisfied by different input UTXOs,
enabling output-centric covenants. However, it also means a second deposit to the same
CTV address creates a UTXO that cannot satisfy the template---the committed output
amounts are derived from the original deposit, and a different deposit amount produces
a permanently unspendable UTXO.

The CTV vault uses a two-step transaction tree. The vault UTXO is encumbered by a bare
CTV script committing to the unvault transaction. The unvault transaction creates a
P2WSH output with two spending paths: a hot-key withdrawal after a CSV delay (the
normal path), and a CTV-committed cold sweep (the recovery path). Both paths include
an anchor output for CPFP fee bumping. The entire tree is deterministic---it can be
regenerated from key material and the original deposit amount without storing presigned
transactions.

\begin{figure}[t]
\centering
\resizebox{0.95\textwidth}{!}{%
\begin{tikzpicture}[
  utxo/.style={rectangle, draw, rounded corners=2pt, minimum width=2.8cm,
               minimum height=1.0cm, align=center, font=\footnotesize,
               fill=blue!5, draw=blue!60!black, thick},
  script/.style={rectangle, draw, rounded corners=1pt, minimum width=2.8cm,
                 minimum height=0.5cm, align=center, font=\scriptsize\ttfamily,
                 fill=yellow!20, draw=yellow!60!black},
  leaf/.style={rectangle, draw, rounded corners=2pt, minimum width=2.6cm,
               minimum height=0.9cm, align=center, font=\scriptsize, thick},
  anchorstyle/.style={rectangle, draw, rounded corners=2pt, minimum width=2.2cm,
                      minimum height=0.7cm, align=center, font=\scriptsize,
                      fill=orange!10, draw=orange!60!black, thick},
  note/.style={font=\scriptsize, align=center, text=gray!70!black},
  phase/.style={font=\scriptsize\bfseries, text=gray!60!black},
  edge/.style={draw, ->, >=stealth, thick},
  hashbox/.style={rectangle, draw, dashed, rounded corners=2pt,
                  align=left, font=\tiny\ttfamily,
                  fill=gray!3, inner sep=5pt},
]
  % Phase labels (top row)
  \node[phase] (pd) at (0, 2.2) {Deposit};
  \node[phase] (pt) at (4.2, 2.2) {Trigger};
  \node[phase] (pu) at (8.4, 2.2) {Unvault output};
  \node[phase] (ps) at (12.6, 2.2) {Spending paths};

  % Deposit: vault UTXO
  \node[utxo] (vault) at (0, 0.9) {Vault UTXO\\ \footnotesize 0.5 BTC};
  \node[script, below=0.05cm of vault] (vaultscript) {<H> OP\_CTV};

  % Trigger: unvault tx
  \node[utxo] (unvault) at (4.2, 0.9) {Unvault TX\\ \footnotesize template = H};

  % Unvault output: P2WSH with 2 leaves
  \node[utxo] (unvaultout) at (8.4, 0.9) {Unvault UTXO\\ \footnotesize 0.5 BTC − fees};
  \node[script, below=0.05cm of unvaultout] {P2WSH, 2 leaves};

  % Anchor output (separate row below the tx)
  \node[anchorstyle] (anchor) at (8.4, -1.6) {Anchor\\ \footnotesize 3{,}300 sats};

  % Spending paths (right column)
  \node[leaf, fill=green!10, draw=green!60!black] (withdraw) at (12.6, 1.5)
        {Withdraw\\ \scriptsize hot + CSV delay};
  \node[leaf, fill=red!10, draw=red!60!black] (cold) at (12.6, 0.1)
        {Cold Sweep\\ \scriptsize CTV-committed};

  % Arrows: vault -> unvault (trigger)
  \draw[edge] (vault.east) -- (unvault.west);
  \node[note, above=0.0cm of {(2.1,1.55)}] {spend = template $H$};

  % unvault -> unvaultout
  \draw[edge] (unvault.east) -- (unvaultout.west);

  % unvault -> anchor (dashed, fee-bump branch)
  \draw[edge, dashed, orange!70!black] (unvault.south) -- ++(0,-1.0) -| (anchor.north)
        node[pos=0.75, left, note, text=orange!70!black] {CPFP};

  % unvaultout -> two leaves
  \draw[edge, green!60!black] (unvaultout.east) -- ++(0.6,0) |- (withdraw.west);
  \draw[edge, red!60!black]   (unvaultout.east) -- ++(0.6,0) |- (cold.west);

  % Template hash note (placed low-left, out of other elements' way)
  \node[hashbox] (hashdef) at (0, -2.5)
        {\textbf{Template hash $H$ commits to:}\\
         nVersion, nLockTime,\\
         scriptSigs hash, input count,\\
         nSequence hash, output count,\\
         \textbf{outputs hash}, input index};
  \draw[->, dashed, gray!60] (vaultscript.south) -- (hashdef.north);

\end{tikzpicture}}
\caption{CTV vault mechanism. The vault UTXO is encumbered by a bare CTV hash $H$
that commits to the unvault transaction's entire output structure; any spend
whose template hash differs from $H$ fails. The unvault UTXO is a P2WSH script
with two Taproot leaves: a hot-key withdrawal gated by \texttt{OP\_CSV}, and a
CTV-committed cold sweep. A 3{,}300-sat anchor output enables CPFP fee bumping.}
\label{fig:ctv_mechanism}
\end{figure}

% ======================================================================
\section{OP\_CHECKCONTRACTVERIFY (CCV)} \label{se:ccv}

CCV~\cite{BIP443} verifies relationships between inputs, outputs, and
Taproot~\cite{BIP341, BIP342} contract structure. The opcode takes five arguments from the stack: optional state data, an
output or input index, an expected internal key, an expected taptree root, and a
\emph{mode} byte that controls which checks are performed. Mode~0 verifies the output's
scriptPubKey matches the expected Taproot structure. Mode~1 checks amounts. Mode~2
checks both. Mode~$-1$ verifies the current input rather than an output.

CCV operates at the per-output granularity: each invocation constrains a single output
(or input) by index. This is fundamentally different from CTV's all-outputs-at-once
commitment. The per-output model enables two capabilities absent in CTV: partial
withdrawal via \texttt{trigger\_and\_revault} (splitting a vault into a withdrawal
portion and a re-funded vault in one transaction) and batched triggers (combining
multiple vault UTXOs into a single trigger transaction).

The CCV vault combines CCV for state transitions with CTV for the final withdrawal
template. The trigger step stores a CTV hash into the contract state. The withdraw
path enforces a CSV timelock and executes \texttt{OP\_CHECKTEMPLATEVERIFY} against the
stored hash. Recovery is keyless---anyone can invoke the recover clause, which sends
funds to the pre-committed recovery address without requiring a signature.

An important design property: undefined CCV mode values (any value outside
$\{-1, 0, 1, 2\}$) trigger \texttt{OP\_SUCCESS}, making the entire script succeed
unconditionally. This is deliberate forward-compatibility behavior specified in
BIP-443 to enable future soft-fork extensions, but it creates a developer footgun:
a mode byte error silently disables all covenant enforcement.

\begin{figure}[t]
\centering
\resizebox{0.95\textwidth}{!}{%
\begin{tikzpicture}[
  utxo/.style={rectangle, draw, rounded corners=2pt, minimum width=2.7cm,
               minimum height=0.95cm, align=center, font=\footnotesize,
               fill=blue!5, draw=blue!60!black, thick},
  leaf/.style={rectangle, draw, rounded corners=1pt, minimum width=2.7cm,
               minimum height=0.5cm, align=center, font=\scriptsize\ttfamily,
               fill=yellow!20, draw=yellow!60!black},
  recleaf/.style={rectangle, draw, rounded corners=1pt, minimum width=2.7cm,
                  minimum height=0.5cm, align=center, font=\scriptsize\ttfamily,
                  fill=red!12, draw=red!60!black},
  opcode/.style={rectangle, draw, rounded corners=2pt, minimum width=3.1cm,
                 minimum height=1.4cm, align=center, font=\scriptsize,
                 fill=purple!8, draw=purple!60!black, thick},
  modebox/.style={rectangle, draw, rounded corners=2pt, dashed,
                  align=left, font=\tiny\ttfamily,
                  fill=gray!3, inner sep=5pt},
  phase/.style={font=\scriptsize\bfseries, text=gray!60!black, align=center},
  note/.style={font=\scriptsize, align=center, text=gray!70!black},
  edge/.style={draw, ->, >=stealth, thick},
]
  % Phase labels across top
  \node[phase] at (0, 3.5) {Deposit (2 leaves)};
  \node[phase] at (4.5, 3.5) {Per-output check};
  \node[phase] at (9.0, 3.5) {Trigger outputs};

  % Deposit: vault with TWO leaves (trigger + recover)
  \node[utxo] (vault) at (0, 2.3) {Vault UTXO\\ \footnotesize 0.5 BTC};
  \node[leaf, below=0.1cm of vault] (triggerleaf) {trigger leaf: OP\_CCV};
  \node[recleaf, below=0.1cm of triggerleaf] (recoverleaf) {recover leaf: keyless};

  % CCV opcode detail
  \node[opcode] (ccv) at (4.5, 2.0)
       {\textbf{OP\_CCV}\\
        \ttfamily mode $\in$ \{−1,0,1,2\}\\
        \ttfamily <idx, pk,\\ \ttfamily taptree, data>};

  % Split outputs (trigger path)
  \node[utxo] (revault) at (9.0, 2.8) {Revault UTXO\\ \footnotesize 0.3 BTC};
  \node[leaf, below=0.05cm of revault] {CCV (recursive)};

  \node[utxo, fill=green!10, draw=green!60!black] (wdraw) at (9.0, 1.1)
       {Withdraw UTXO\\ \footnotesize 0.2 BTC};
  \node[leaf, below=0.05cm of wdraw] {CTV + CSV};

  % Recovery output (separate lane, below)
  \node[utxo, fill=red!10, draw=red!60!black] (rec) at (9.0, -1.8)
       {Recovery UTXO\\ \footnotesize \textbf{keyless}};
  \node[leaf, below=0.05cm of rec] {fixed address};

  % ===== Trigger path: trigger leaf -> OP_CCV -> [revault, withdraw] =====
  \draw[edge] (triggerleaf.east) -- ++(0.3,0) |- (ccv.west)
        node[pos=0.6, above, note] {trigger};
  \draw[edge] (ccv.east) -- ++(0.5,0) |- (revault.west);
  \draw[edge] (ccv.east) -- ++(0.5,0) |- (wdraw.west);

  % ===== Recovery path: recover leaf -> Recovery UTXO (DIRECT, bypasses OP_CCV) =====
  \draw[edge, red!70!black, rounded corners=6pt]
       (recoverleaf.east) -- ++(0.3,0) -- ++(0,-1.3) -| (rec.south);
  \node[note, text=red!70!black] at (5.2, -1.6) {any party may invoke};

  % Mode legend — placed above-right of OP_CCV to avoid crossing the recovery path
  \node[modebox] (modes) at (5.2, 5.3)
       {\textbf{Mode semantics:}\\
        mode\,0: check spk of output[idx]\\
        mode\,1: check amount[idx]\\
        mode\,2: check both\\
        mode\,-1: check this input};
  \draw[->, dashed, purple!60!black] (modes.south) -- (ccv.north);

\end{tikzpicture}}
\caption{CCV vault mechanism. The vault UTXO is a Taproot with two leaves:
a \emph{trigger} leaf containing \texttt{OP\_CCV}, and a keyless \emph{recover}
leaf. Spending the trigger leaf invokes the per-output
\texttt{OP\_CHECKCONTRACTVERIFY} check, whose mode byte selects whether to
verify the scriptPubKey, the amount, both, or the current input. Taptree
preservation enables recursion, supporting partial withdrawal via
\texttt{trigger\_and\_revault}. Spending the recover leaf is permissionless
and sends funds to the pre-committed address without a signature.}
\label{fig:ccv_mechanism}
\end{figure}

% ======================================================================
\section{OP\_VAULT (BIP-345)} \label{se:opvault}

OP\_VAULT~\cite{BIP345} provides dedicated vault opcodes: \texttt{OP\_VAULT} controls
withdrawal initiation, and \texttt{OP\_VAULT\_RECOVER} enables authorized recovery.
Unlike CCV's generic contract verification, these opcodes encode vault-specific semantics
directly. Amount preservation (ensuring outputs account for all input value), automatic
revault (excess value beyond the withdrawal amount is re-locked), and recovery
authorization (requiring a \texttt{recoveryauth} key signature) are enforced at the
opcode level.

The OP\_VAULT design uses a three-key model: a \emph{trigger key} (derived from an xpub)
initiates withdrawals, a \emph{recoveryauth key} authorizes emergency recovery, and a
\emph{recovery address} (embedded in the vault's Taproot internal key tweak) is the
pre-committed destination for recovered funds. This three-key separation is unique among
the designs studied: CTV and CAT+CSFS use two keys (hot + cold), and CCV uses one
(trigger key; recovery is keyless).

A structural consequence of OP\_VAULT's design is the \emph{fee-input requirement}. The
opcode's value preservation rule requires that all input value be accounted for in
outputs. Since transaction fees are paid from input value, the vault UTXO alone cannot
cover fees without violating value preservation. The reference implementation solves this
by adding a separate fee-wallet UTXO as a second input to every trigger and recovery
transaction, adding 80--90\vB{} of overhead per transaction.

BIP-345 was withdrawn in early 2025 in favor of BIP-443~\cite{OB25}. We include it in
this comparison because (a)~it represents a distinct design philosophy (purpose-built
vs.\ generic covenants), (b)~the economic comparison with CCV quantifies the cost of
that design choice (36\% lifecycle overhead), and (c)~it serves as a historical reference
for the OP\_VAULT-to-CCV transition.

\begin{figure}[t]
\centering
\resizebox{0.95\textwidth}{!}{%
\begin{tikzpicture}[
  utxo/.style={rectangle, draw, rounded corners=2pt, minimum width=2.8cm,
               minimum height=0.95cm, align=center, font=\footnotesize,
               fill=blue!5, draw=blue!60!black, thick},
  fee/.style={rectangle, draw, rounded corners=2pt, minimum width=2.6cm,
              minimum height=0.75cm, align=center, font=\scriptsize,
              fill=orange!15, draw=orange!70!black, thick},
  leaf/.style={rectangle, draw, rounded corners=1pt, minimum width=2.8cm,
               minimum height=0.5cm, align=center, font=\scriptsize\ttfamily,
               fill=yellow!20, draw=yellow!60!black},
  keynode/.style={circle, draw, minimum size=0.95cm, font=\tiny, align=center,
              inner sep=0pt},
  phase/.style={font=\scriptsize\bfseries, text=gray!60!black, align=center},
  note/.style={font=\scriptsize, align=center, text=gray!70!black},
  edge/.style={draw, ->, >=stealth, thick},
]
  % Column x-coordinates
  \def\xkey{-0.2}
  \def\xvault{3.1}
  \def\xtx{7.0}
  \def\xout{11.0}

  % Phase labels
  \node[phase] at (\xkey, 3.4)   {Keys};
  \node[phase] at (\xvault, 3.4) {Vault UTXO};
  \node[phase] at (\xtx, 3.4)    {Transactions};
  \node[phase] at (\xout, 3.4)   {Outputs};

  % Three keys (left column)
  \node[keynode, fill=green!15] (tk) at (\xkey, 2.3) {trigger\\ key};
  \node[keynode, fill=red!15]   (ak) at (\xkey, 0.3) {recovery-\\ auth};
  \node[keynode, fill=blue!15]  (rk) at (\xkey, -1.7) {recovery\\ address};

  % Vault UTXO + two leaves
  \node[utxo] (vault) at (\xvault, 1.3) {Vault UTXO\\ \footnotesize 0.5 BTC};
  \node[leaf, below=0.05cm of vault] (tleaf) {OP\_VAULT};
  \node[leaf, below=0.1cm of tleaf] (rleaf) {OP\_VAULT\_RECOVER};

  % Fee wallet (above trigger)
  \node[fee] (feewallet) at (\xtx, 3.9) {Fee-wallet UTXO};
  \node[note, below=0.0cm of feewallet, yshift=-0.4cm, fill=orange!5,
        rounded corners=1pt, inner sep=2pt, font=\tiny\itshape, text=orange!70!black]
        {required 2nd input};

  % Trigger TX
  \node[utxo] (trigger) at (\xtx, 1.5) {Trigger TX};

  % Recovery TX (below)
  \node[utxo, fill=red!8, draw=red!70!black] (recovery) at (\xtx, -1.7) {Recovery TX};

  % Outputs
  \node[utxo, fill=green!10, draw=green!60!black] (wdraw) at (\xout, 2.2)
       {Withdrawal\\ \footnotesize CSV-delayed};
  \node[utxo] (revault2) at (\xout, 0.6) {Revault UTXO\\ \footnotesize auto-created};
  \node[utxo, fill=red!10, draw=red!60!black] (recdest) at (\xout, -1.7)
       {Recovery Addr\\ \footnotesize pre-committed};

  % Key authorization edges (dashed, short)
  \draw[edge, dashed, green!60!black] (tk.east) -- ++(0.4,0) |- (tleaf.west);
  \draw[edge, dashed, red!70!black]   (ak.east) -- ++(0.4,0) |- (rleaf.west);

  % Recovery-address key: routed below everything, up into recdest from below
  \draw[edge, dashed, blue!70!black]
       (rk.west) -- ++(-0.5,0) -- ++(0,-1.2)
       -- (\xout, -3.2) -- (recdest.south);

  % Vault -> Trigger
  \draw[edge] (tleaf.east) -- ++(0.4,0) |- (trigger.west);

  % Fee wallet -> Trigger (primary, dashed orange)
  \draw[edge, dashed, orange!70!black] (feewallet.south) -- (trigger.north);

  % Vault recover -> Recovery TX
  \draw[edge, red!60!black] (rleaf.east) -- ++(0.4,0) |- (recovery.west);

  % Fee wallet -> Recovery TX: route down the LEFT side so it doesn't cross Recovery Addr
  \draw[edge, dashed, orange!70!black]
       (feewallet.west) -- ++(-0.4,0) |- (recovery.west);

  % Trigger -> outputs
  \draw[edge, green!60!black] (trigger.east) -- ++(0.3,0) |- (wdraw.west);
  \draw[edge] (trigger.east) -- ++(0.3,0) |- (revault2.west);

  % Recovery -> dest
  \draw[edge, red!60!black] (recovery.east) -- (recdest.west);

\end{tikzpicture}}
\caption{OP\_VAULT mechanism. The opcode pair \texttt{OP\_VAULT} and
\texttt{OP\_VAULT\_RECOVER} encodes vault-specific semantics (amount preservation,
auto-revault, authorized recovery) directly. Three keys are separated: the
\emph{trigger key} initiates withdrawals, the \emph{recoveryauth key} authorizes
recovery, and the \emph{recovery address} (embedded in the Taproot internal-key
tweak) is the pre-committed destination. Value preservation requires a fee-wallet
input on every trigger and recovery transaction (${\sim}80$--$90\vB{}$ overhead).}
\label{fig:opvault_mechanism}
\end{figure}

% ======================================================================
\section{CAT+CSFS (BIP-347 + BIP-348)} \label{se:catcsfs_bg}

CAT+CSFS achieves transaction introspection by composing two general-purpose
opcodes~\cite{Poe21, BIP347, BIP348}. \texttt{OP\_CAT} (BIP-347) concatenates the top
two stack elements. \texttt{OP\_CHECKSIGFROMSTACK} (BIP-348) verifies a Schnorr
signature against an arbitrary stack-provided message rather than the transaction's
sighash. The introspection pattern works as follows:

\begin{enumerate}
\item The spender provides the sighash preimage fields (transaction version, hash
  prevouts, hash amounts, hash script pubkeys, hash sequences, hash outputs, spend
  type, and other BIP-342 fields) as individual witness elements.
\item The script uses \texttt{OP\_CAT} to assemble these fields into the complete
  sighash preimage on the stack.
\item \texttt{OP\_CHECKSIGFROMSTACK} verifies a Schnorr signature against this
  assembled preimage using the hot key.
\item \texttt{OP\_CHECKSIG} verifies the \emph{same signature} against the real
  transaction sighash computed by the Bitcoin consensus layer.
\end{enumerate}

If both checks pass with the same key and signature, the witness-provided preimage
must be the actual transaction's preimage (otherwise the signature could not satisfy
both). The script can embed constants (such as \texttt{sha\_single\_output}, the SHA-256
hash of the expected output) that the spender cannot alter, effectively constraining
the transaction's outputs through the signature binding.

The CAT+CSFS vault uses \texttt{SIGHASH\_SINGLE|ANYONECANPAY} (0x83), a sighash
mode defined in BIP-342~\cite{BIP342} that commits to one input and one output.
This allows external fee-paying inputs to be attached without breaking the
covenant---a more flexible fee model than CTV's anchor-based CPFP.

Two properties of this design have direct security implications. First, the withdrawal
destination is fixed at vault creation time (the \texttt{sha\_single\_output} constant
is embedded in the script), making it the most rigid destination constraint of any
design studied. Second, the recovery leaf uses bare \texttt{cold\_pk OP\_CHECKSIG}
with no introspection, meaning cold key compromise grants the attacker unrestricted
control over recovery---the weakest recovery path among the four Bitcoin designs.

\begin{figure}[t]
\centering
\resizebox{0.97\textwidth}{!}{%
\begin{tikzpicture}[
  witness/.style={rectangle, draw, minimum width=2.2cm, minimum height=0.45cm,
                  font=\tiny\ttfamily, align=center, fill=gray!10},
  opbox/.style={rectangle, draw, rounded corners=2pt, minimum width=3.0cm,
                minimum height=0.9cm, align=center, font=\scriptsize,
                fill=purple!10, draw=purple!60!black, thick},
  sigbox/.style={rectangle, draw, rounded corners=1pt, minimum width=2.8cm,
                 minimum height=0.6cm, align=center, font=\scriptsize,
                 fill=cyan!10, draw=cyan!60!black},
  leafcontainer/.style={rectangle, draw, rounded corners=3pt, thick,
                        dash pattern=on 3pt off 2pt},
  bad/.style={rectangle, draw, rounded corners=2pt, minimum width=3.0cm,
              minimum height=0.8cm, align=center, font=\scriptsize,
              fill=red!8, draw=red!70!black, thick},
  note/.style={font=\scriptsize, align=center, text=gray!70!black},
  phase/.style={font=\scriptsize\bfseries, text=gray!60!black, align=center},
  edge/.style={draw, ->, >=stealth, thick},
]
  % ========= HOT LEAF (top container) =========
  \node[leafcontainer, fill=green!2, draw=green!50!black, minimum width=15.5cm,
        minimum height=5.2cm] (hotleaf) at (7.2, 0.7) {};
  \node[font=\small\bfseries, text=green!60!black, anchor=west]
       at ($(hotleaf.north west)+(0.2,-0.25)$) {Hot-key leaf (dual verification)};

  % Phase labels along hot leaf
  \node[phase] at (0.5, 2.7) {Witness fields};
  \node[phase] at (4.9, 2.7) {Assemble};
  \node[phase] at (8.6, 2.7) {Preimage};
  \node[phase] at (12.2, 2.7) {Dual verification};

  % Witness fields
  \node[witness] (w1) at (0.5, 2.0) {nVersion};
  \node[witness, below=0.02cm of w1] (w2) {hashPrevouts};
  \node[witness, below=0.02cm of w2] (w3) {hashAmounts};
  \node[witness, below=0.02cm of w3] (w4) {hashSPKs};
  \node[witness, below=0.02cm of w4] (w5) {hashSequences};
  \node[witness, below=0.02cm of w5, fill=red!15, draw=red!60!black]
       (w6) {\textbf{hashOutputs}};
  \node[witness, below=0.02cm of w6] (w7) {spendType, etc.};

  % OP_CAT
  \node[opbox] (cat) at (4.9, 0.5) {\textbf{OP\_CAT}\\ concatenate};

  % Preimage
  \node[sigbox] (preimage) at (8.6, 0.5) {sighash preimage\\ (assembled)};

  % OP_CSFS (upper)
  \node[opbox] (csfs) at (12.2, 1.8) {\textbf{OP\_CSFS}\\ \scriptsize sig, hot\_pk,
        assembled};
  % OP_CHECKSIG (lower)
  \node[opbox, fill=teal!10, draw=teal!70!black] (csig) at (12.2, -0.8)
       {\textbf{OP\_CHECKSIG}\\ \scriptsize sig, hot\_pk, \textbf{real sighash}};

  % Binding outcome (right of hot leaf, same level)
  \node[sigbox, fill=green!12, draw=green!60!black, minimum width=2.6cm,
        minimum height=0.9cm, font=\scriptsize\bfseries] (bind)
       at (15.9, 0.5) {Output\\ bound};

  % Edges inside hot leaf
  \foreach \i in {1,2,3,4,5,7} {
    \draw[edge, gray!70] (w\i.east) -- ++(0.2,0) |- (cat.west);
  }
  \draw[edge, red!60!black, thick] (w6.east) -- ++(0.2,0) |- (cat.west);

  \draw[edge] (cat) -- (preimage);
  \draw[edge, purple!60!black] (preimage.east) -- ++(0.25,0) |- (csfs.west);
  \draw[edge, teal!70!black] (preimage.east) -- ++(0.25,0) |- (csig.west);
  % Both verifications converge on Output bound's west face via L-shapes
  \draw[edge, purple!60!black] (csfs.east) -| (bind.north west);
  \draw[edge, teal!70!black]   (csig.east) -| (bind.south west);
  \node[font=\tiny\itshape, text=gray!70!black] at (14.55, 0.5) {$\wedge$};

  % ========= RECOVERY LEAF (bottom container, contrast) =========
  \node[leafcontainer, fill=red!2, draw=red!60!black, minimum width=15.5cm,
        minimum height=1.5cm] (coldleaf) at (7.2, -3.0) {};
  \node[font=\small\bfseries, text=red!70!black, anchor=west]
       at ($(coldleaf.north west)+(0.2,-0.25)$) {Recovery leaf (no binding)};

  \node[bad] (recoverbox) at (4.9, -3.2) {cold\_pk OP\_CHECKSIG};
  \node[note, text=red!70!black, align=left, font=\scriptsize\itshape]
       at (10.0, -3.2) {no CAT, no CSFS, no \texttt{hashOutputs} binding\\
                         → attacker chooses destination};

\end{tikzpicture}}
\caption{CAT+CSFS vault mechanism. The Taproot script has two leaves.
\textbf{Hot-key leaf (top):} the spender provides sighash preimage fields;
\texttt{OP\_CAT} concatenates them; \texttt{OP\_CSFS} verifies the Schnorr
signature against this stack-assembled preimage; \texttt{OP\_CHECKSIG} verifies
the \emph{same} signature against the real transaction sighash. Both succeed
only if the two preimages agree, binding outputs through \texttt{hashOutputs}.
\textbf{Recovery leaf (bottom):} bare \texttt{cold\_pk OP\_CHECKSIG}---no CAT,
no CSFS, no output binding.}
\label{fig:catcsfs_mechanism}
\end{figure}

% ======================================================================
\section{Simplicity on Elements: An Alternative Script Primitive} \label{se:simplicity_bg}

The four proposals above extend Bitcoin Script by introducing new
opcodes or composing existing ones under soft-fork consensus change.
Simplicity takes a different path: it is an \emph{alternative script
primitive}, a typed combinator calculus with denotational semantics
verified in the Coq proof assistant, designed by O'Connor~\cite{OC17}.
Programs are directed acyclic graphs (DAGs) of combinators. The
language runs on Elements/Liquid~\cite{DPNS16}, a federated Bitcoin
sidechain. Both Bitcoin Script and Simplicity are scripting languages;
they differ in paradigm (stack-based vs.\ typed functional), formal
foundation (informal specification vs.\ mechanised semantics), and
operating substrate (Bitcoin L1 proof-of-work vs.\ Elements federated
consensus).

Execution uses \emph{jets}: when the validator encounters a known combinator subtree
(identified by its canonical Merkle root), it executes a native C implementation instead
of interpreting the combinators. Jets for transaction introspection
(\texttt{jet::outputs\_hash()}, \texttt{jet::sig\_all\_hash()}), timelocks
(\texttt{jet::check\_lock\_distance()}), and cryptographic primitives
(\texttt{jet::bip\_0340\_verify()}~\cite{BIP340}) provide the building blocks for
covenant enforcement.

The Simplicity vault enforces covenants via \texttt{jet::outputs\_hash()}, which
computes a Merkle hash over all outputs of the spending transaction. The expected hash
is embedded as a compile-time parameter in the Simplicity program. At spend time, the
jet computes the hash from the actual transaction outputs; if the values differ, the
program fails at the consensus level. This is structurally analogous to CTV's template
hash but operates over Elements' richer output model (explicit asset IDs, confidential
values, range proofs).

The Simplicity vault runs on Elements, a federated sidechain where block production is
controlled by known functionaries rather than proof-of-work miners. This fundamentally
changes the trust model: functionaries could censor attacker transactions or collude
with an attacker, making the threat models (TM1--TM11) applicable only under the
assumption of honest federation behavior. Measurements from Elements are presented
separately from the four Bitcoin covenant measurements and are not directly comparable
due to differences in transaction serialization (explicit fee outputs, confidential
transaction fields add approximately 30--40 bytes per output).

\begin{figure}[t]
\centering
\resizebox{0.95\textwidth}{!}{%
\begin{tikzpicture}[
  prog/.style={rectangle, draw, rounded corners=3pt, minimum width=4.6cm,
               minimum height=4.0cm, align=center, font=\scriptsize,
               fill=purple!4, draw=purple!60!black, thick},
  combinator/.style={circle, draw, minimum size=0.65cm, font=\tiny,
                     inner sep=0pt, fill=white},
  jetnode/.style={rectangle, draw, rounded corners=2pt, minimum width=2.0cm,
                  minimum height=0.55cm, font=\tiny, align=center,
                  fill=orange!20, draw=orange!70!black, thick},
  tx/.style={rectangle, draw, rounded corners=2pt, minimum width=3.2cm,
             minimum height=1.0cm, align=center, font=\footnotesize,
             fill=blue!5, draw=blue!60!black, thick},
  jetbig/.style={rectangle, draw, rounded corners=2pt, minimum width=3.6cm,
                 minimum height=0.9cm, font=\scriptsize, align=center,
                 fill=orange!20, draw=orange!70!black, thick},
  paramnodestyle/.style={rectangle, draw, rounded corners=2pt, minimum width=3.0cm,
                         minimum height=0.7cm, align=center, font=\scriptsize,
                         fill=cyan!10, draw=cyan!60!black, thick},
  checkbox/.style={diamond, draw, aspect=1.8, minimum width=2.4cm,
                   font=\scriptsize, align=center, thick,
                   fill=yellow!15, draw=yellow!70!black},
  outcome/.style={rectangle, draw, rounded corners=2pt, minimum width=2.4cm,
                  minimum height=0.9cm, align=center, font=\scriptsize\bfseries,
                  thick},
  phase/.style={font=\scriptsize\bfseries, text=gray!60!black, align=center},
  note/.style={font=\scriptsize, align=center, text=gray!70!black},
  edge/.style={draw, ->, >=stealth, thick},
]
  % ===== Phase headers (top, above all content) =====
  \node[phase] at (0, 3.9)    {Compile time};
  \node[phase] at (5.2, 3.9)  {Spend time};
  \node[phase] at (9.7, 3.9)  {Consensus};
  \node[phase] at (13.5, 3.9) {Outcome};

  % ===== Compile-time column: program with embedded H =====
  \node[prog] (program) at (0, 1.2) {};
  \node[font=\scriptsize\itshape, anchor=north]
       at ($(program.north)+(0,-0.15)$) {Simplicity program};

  % DAG mini-diagram (upper part of program)
  \node[combinator] (c1) at (-1.5, 2.3) {pair};
  \node[combinator] (c2) at (-1.5, 1.3) {iden};
  \node[combinator] (c3) at (-0.4, 1.3) {comp};
  \node[jetnode] (j1) at (1.1, 2.3) {\ttfamily outputs\_hash};
  \node[jetnode] (j2) at (1.1, 1.3) {\ttfamily sig\_all\_hash};
  \draw[->, thin, gray] (c1) -- (j1);
  \draw[->, thin, gray] (c3) -- (j2);
  \draw[->, thin, gray] (c2) -- (c3);
  \draw[->, thin, gray] (c1) -- (c3);

  % H parameter (lower part of program box)
  \node[paramnodestyle] (paramnode) at (0, 0.0) {$H$ = expected outputs hash};

  % ===== Spend-time column =====
  \node[tx] (spendtx) at (5.2, 2.6) {Spending TX\\ \footnotesize outputs: [out1, out2]};
  \node[jetbig] (jeteval) at (5.2, 0.5) {\ttfamily jet::outputs\_hash()};
  \node[note, below=0.1cm of jeteval, align=center, text width=4cm]
       {computes hash at spend time};

  % ===== Consensus column =====
  \node[checkbox] (check) at (9.7, 1.5) {$H =$\\ computed?};

  % ===== Outcome column =====
  \node[outcome, fill=green!15, draw=green!70!black] (ok) at (13.5, 2.5) {Accept};
  \node[outcome, fill=red!15, draw=red!70!black]   (fail) at (13.5, 0.3)
       {Consensus\\ fail};

  % ===== Edges =====
  % Spending TX -> jet (labeled "outputs")
  \draw[edge] (spendtx.south) -- (jeteval.north)
        node[midway, right=1pt, note] {outputs};

  % jet -> check (upper entry)
  \draw[edge] (jeteval.east) -- ++(0.5,0) |- (check.south)
        node[pos=0.9, below, note] {computed hash};

  % H -> check: route from paramnode top-right, rise to check level, into check.west
  \draw[edge, cyan!70!black]
       (paramnode.east) -- ++(0.6,0) -- (3.3, 1.5) -- (check.west)
       node[pos=0.95, above, note, text=cyan!70!black] {$H$};

  % Program "invokes" jet at spend time (dashed)
  \draw[edge, dashed, purple!60!black]
       (program.east |- jeteval.west) -- (jeteval.west)
       node[pos=0.4, above, note, text=purple!60!black] {invokes};

  % check -> outcomes
  \draw[edge, green!60!black] (check.east) -- ++(0.9,0) |- (ok.west)
        node[pos=0.8, above, note] {yes};
  \draw[edge, red!60!black]   (check.east) -- ++(0.9,0) |- (fail.west)
        node[pos=0.8, below, note] {no};

\end{tikzpicture}}
\caption{Simplicity vault mechanism on Elements. At compile time, the program
(a DAG of combinators) embeds an expected outputs-hash~$H$ as a parameter.
At spend time, \texttt{jet::outputs\_hash()} computes the Merkle hash of the
spending transaction's outputs. If the computed hash matches~$H$, the program
accepts; otherwise the program fails at the consensus level. Both trigger and
recovery leaves enforce this jet---unlike CAT+CSFS, where recovery is
unconstrained.}
\label{fig:simplicity_mechanism}
\end{figure}

\section{The Design-Space Axes}
\label{se:design_space_axes}

The four Bitcoin covenant proposals differ along seven axes.
Three carry over from the formal decomposition (\cite{SHMB20,Har24}
and this work) and appear with Greek-letter symbols $(f, a, g, b)$ in
proofs (Chapter~\ref{ch:security}); the remaining four are introduced
in this thesis and are used as descriptive labels. Every attack class
in Chapter~\ref{ch:security} derives from a specific axis-value;
every design's vulnerability profile is a lookup against
Table~\ref{tab:design_space_position}. We define the axes here, in
Background, so that Chapter~\ref{ch:security}'s analysis reads as
mechanical derivation rather than assertion.

\subsection{Axis Definitions}
\label{ss:axis_defs}

\begin{description}
\item[$\axauthf{}$ --- Recovery authorisation.] Who may initiate the
  recovery path? Values: \texttt{keyless}, \texttt{recoveryauth-key},
  \texttt{hot-key}, \texttt{cold-key}. Controls class \classgrief{}
  (permissionless griefing).

\item[$\axrevaultf{}$ --- Revault support.] Does the covenant allow
  partial withdrawal --- splitting one vault UTXO into a withdrawn
  output and a remaining vault UTXO in one transaction? Values:
  \texttt{yes}, \texttt{no}. Controls class \classscale{} (per-UTXO
  recovery-cost scaling).

\item[$\axfeef{}$ --- Fee-inclusion model.] Where does the spending
  transaction's fee come from? Values: \texttt{dynamic} (fee in the
  transaction), \texttt{static-precommit} (fee baked at creation),
  \texttt{out-of-band-anchor} (CPFP via a separate anchor key). Controls
  class \classfee{} (fee-channel pinning).

\item[\axwdbind{} --- Withdrawal-destination binding.] How is the
  withdrawal destination constrained? Values:
  \texttt{template-hash-digest}, \texttt{per-output-amount+script},
  \texttt{dedicated-opcode}, \texttt{signature-dual-bind}. Controls
  class \classhot{} (hot-key theft via destination substitution).

\item[\axrecbind{} --- Recovery-destination binding.] Same question
  for the recovery path. Values: same as \axwdbind{}, plus
  \texttt{plain-signature} (no script-level constraint). Controls
  class \classcold{} (cold-key theft via unconstrained recovery).
  Kept separate from \axwdbind{} because \catcsfs{} occupies different
  positions on the two sub-axes --- strongest on \axwdbind{}, weakest
  on \axrecbind{}. A collapsed binding axis $b$ is used in proofs that
  do not depend on this asymmetry.

\item[\axopsurf{} --- Opcode surface.] Is the covenant expressed by
  a single semantic opcode or a multi-mode polymorphic one with
  \texttt{OP\_SUCCESSx} fallthrough on undefined modes? Values:
  \texttt{single-mode}, \texttt{multi-mode-with-OP\_SUCCESS}. Controls
  class \classmode{} (mode/parameter bypass).

\item[\axintro{} --- Introspection primitive.] How does the script
  read transaction data? Values correspond to the four primitives
  enumerated in Section~\ref{se:introspection}:
  \texttt{digest-equality} (CTV),
  \texttt{merkle-membership-and-amount} (CCV),
  \texttt{dedicated-opcode} (OP\_VAULT), and
  \texttt{sighash-preimage-reconstruction} (CAT+CSFS). Descriptive;
  governs which of the other axes are simultaneously realisable.
\end{description}

For the formal decomposition in Chapter~\ref{ch:security} we retain
$(f, a, g, b)$ as the axes of the design-space lattice used in
Propositions~1 and~2: $f = \axfee$, $a = \axrevault$, $g = \axauth$,
and $b$ is the joint $(\axwdbind, \axrecbind)$ sublattice. Axes
\axopsurf{} and \axintro{} are additional dimensions orthogonal to
the proofs; they are where the developer-footgun (class \classmode{})
and design-implementability questions live.

\subsection{Four-Covenant Position Table}
\label{ss:position_table}

Table~\ref{tab:design_space_position} locates each of the four
Bitcoin proposals on every axis. Every attack-class
susceptibility claim in Chapter~\ref{ch:security} is a mechanical
lookup against this table; every immunity claim is a mechanical
reading of the opposite axis-value. Simplicity on Elements is
treated separately in Section~\ref{se:simplicity_bg} and (at greater
length) in Appendix~C; it occupies the strict-best position on
\axauth{}, \axrevault{}, \axwdbind{}, \axrecbind{}, and \axintro{}
simultaneously, and no Bitcoin proposal does.

\begin{table}[t]
\centering
\caption{The four Bitcoin covenant proposals, positioned on
the seven axes of Section~\ref{ss:axis_defs}. Every attack-class
susceptibility and immunity claim in Chapter~\ref{ch:security} is a
lookup against this table; numerical thresholds are measured in
Chapter~\ref{ch:results}.}
\label{tab:design_space_position}
\footnotesize
\setlength{\tabcolsep}{2.2pt}
\renewcommand{\arraystretch}{1.12}
\begin{tabular}{@{}lllll@{}}
\toprule
Axis                                    & \opctv{}               & \opccv{}                     & \opvault{}           & \catcsfs{}                      \\
\midrule
\axauth{} (recovery auth)               & hot-key                & keyless                      & recoveryauth-key     & cold-key                        \\
\axrevault{} (revault)                  & no                     & yes                          & yes                  & no                              \\
\axfee{} (fee model)                    & out-of-band-anchor     & dynamic                      & dynamic              & dynamic                         \\
\axwdbind{} (withdrawal binding)        & template-hash-digest   & per-output-amount+script     & dedicated-opcode     & signature-dual-bind             \\
\axrecbind{} (recovery binding)         & template-hash-digest   & per-output-amount+script     & dedicated-opcode     & plain-signature                 \\
\axopsurf{} (opcode surface)            & single-mode            & multi-mode-with-OP\_SUCCESS  & single-mode          & single-mode                     \\
\axintro{} (introspection primitive)    & digest-equality        & merkle-membership-and-amount & dedicated-opcode     & sighash-preimage-reconstruction \\
\bottomrule
\end{tabular}
\end{table}

\subsection{Structural Dependencies Between Axes}
\label{ss:axis_deps}

Two spec-level dependencies between axes constrain the realisable
combinations and are used throughout Chapter~\ref{ch:security}:

\begin{enumerate}
\item \axfee{} = \texttt{out-of-band-anchor} implies \axrevault{} =
  \texttt{no}. A CPFP-anchor fee model relies on a template hash
  that commits to the whole output set; splitting the vault into
  $[\text{withdrawn}] + [\text{remaining}]$ would require the
  template to anticipate the split amount, which the user cannot
  know at vault creation. CTV's position on these two axes is not
  coincidental; it is forced.

\item \axintro{} = \texttt{sighash-preimage-reconstruction} is the
  only known introspection primitive that admits \axwdbind{} =
  \texttt{signature-dual-bind}. The binding reduces to Schnorr
  message-equality under a single key, which requires reassembling
  the sighash preimage from stack contents and verifying two
  signatures under the same key --- a construction that the
  whole-template-digest, per-output, and dedicated-opcode primitives
  do not support.
\end{enumerate}

These dependencies reduce the effective design space from the naive
$4 \cdot 2 \cdot 3 \cdot 4 \cdot 5 \cdot 2 \cdot 4 = 1920$ cells to
a much smaller set of realisable positions. Every Bitcoin
proposal occupies a realisable position; the thesis contribution is
to characterise every realisable position's attack surface.
